\documentclass[11pt,a4paper]{article}
\usepackage[margin=2.5cm]{geometry}
\usepackage{fontspec}
\setmainfont{Times New Roman}
\usepackage{amsmath,amssymb}
\usepackage{booktabs,tabularx,array}
\usepackage{enumitem}
\usepackage{xcolor}
\usepackage{float}
\usepackage{titlesec}
\usepackage[hidelinks]{hyperref}

\titleformat{\section}{\Large\bfseries\color{black}}{\thesection}{0.6em}{}
\titleformat{\subsection}{\large\bfseries\color{black}}{\thesubsection}{0.6em}{}
\titleformat{\subsubsection}{\normalsize\bfseries\color{black}}{\thesubsubsection}{0.6em}{}

\setlength{\parindent}{0pt}
\setlength{\parskip}{5pt}
\linespread{1.2}
\renewcommand{\arraystretch}{1.15}
\renewcommand{\tabularxcolumn}[1]{>{\raggedright\arraybackslash}p{#1}}
\setlength{\emergencystretch}{4em}

\newcommand{\Cp}{\ensuremath{C_p}}
\newcommand{\CD}{\ensuremath{C_D}}
\newcommand{\CL}{\ensuremath{C_L}}
\newcommand{\ReD}{\ensuremath{\mathit{Re}}}
\newcommand{\thsep}{\ensuremath{\theta_{\text{sep}}}}

\pagestyle{plain}
\begin{document}

\begin{center}
{\LARGE\bfseries Prelab: Pressure Distribution on a Circular Cylinder}\\[4pt]
{\large Engineering Practice 2 -- Aerodynamics 2}\\[8pt]
{\small RODH 115 (ex L50) \quad\textbar\quad $d = 102$\,mm, $b = 460$\,mm, $d/b = 0.2217$
\quad\textbar\quad ethanol manometer $\rho = 789$\,kg\,m$^{-3}$}\\[2pt]
{\small 28 surface tappings \quad\textbar\quad Pitot-static on tubes 35/36 \quad\textbar\quad
data reduction by \texttt{cylinder\_lab.m}}
\end{center}

\vspace{2pt}

%=====================================================================
\section{Purpose of this document}

The handout asks for three things before the session: know the theory, have the reduction code
ready, and be able to answer questions at the bench. This document is the written half of that
preparation. It fixes the definitions and the order of operations now, so that during the lab
nothing has to be derived under time pressure, and it collects the eleven oral-assessment
questions with answers in Section~\ref{sec:oral}.

Everything numerical quoted here (velocities, \Cp{} minima, \CD{} values) is what the code
predicts for this apparatus; the real numbers will be whatever the manometer says.

%=====================================================================
\section{Theory}

\subsection{Pressure coefficient, and how the manometer measures it}

The conventional definition, with $p_\infty$ and $U_\infty$ taken upstream and outside the
blockage-affected region, is
\begin{equation}
\Cp = \frac{p - p_\infty}{\tfrac12 \rho U_\infty^2}.
\end{equation}
Each tapping is connected to one tube of an ethanol manometer, so the pressure at tapping $i$
is $p_i = p_{\text{atm}} + \rho_e g h_i$ with $h_i$ the column height read against the
manifold datum. Taking the static Pitot tube (height $h_s$) as the free-stream reference,
\begin{equation}
p_i - p_\infty = \rho_e g\,(h_i - h_s), \qquad
\tfrac12 \rho U_\infty^2 = q = \rho_e g\,(h_t - h_s),
\end{equation}
and therefore
\begin{equation}\label{eq:cp}
\boxed{\;\ensuremath{C_{p,i}} = \frac{h_i - h_s}{h_t - h_s} = \frac{h_i - h_s}{B}\;}
\end{equation}
The ethanol density, gravity and air density all cancel: \Cp{} is a pure ratio of column heights.
That is worth remembering, because it is the reason a 0.5\,mm reading error is the dominant
uncertainty in the whole experiment, and the reason a low tunnel speed (small $B$) is much
noisier than a high one.

\subsection{The potential-flow answer, and D'Alembert's paradox}

For inviscid, irrotational flow past a cylinder the surface speed is $V_\theta = -2U\sin\theta$,
so Bernoulli gives
\begin{equation}
\Cp(\theta) = 1 - \left(\frac{V_\theta}{U}\right)^2 = 1 - 4\sin^2\theta .
\end{equation}
It follows that $\Cp(0) = +1$ (stagnation), $\Cp(90^\circ) = -3$ (suction peak), and
\begin{equation}
\CD = \int_0^{\pi} \Cp\cos\theta\,\mathrm{d}\theta = 0 ,
\end{equation}
because $\Cp$ is symmetric about $\theta = 90^\circ$ while $\cos\theta$ is antisymmetric. Zero
drag on a body in steady motion is D'Alembert's paradox: the theory has no mechanism to
produce a wake, so the pressure on the front and the back of the cylinder cancel exactly.

\subsection{What viscosity changes: separation and the base-pressure plateau}

With viscosity the no-slip condition creates a boundary layer. From the suction peak
($\theta \approx 90^\circ$ for the ideal case) onwards the surface pressure \emph{rises}
downstream, i.e.\ the boundary layer faces an adverse pressure gradient and loses momentum.
At the point where the wall shear stress vanishes, $\tau_w = 0$, the flow detaches:
\begin{itemize}[leftmargin=1.2em,itemsep=1pt,topsep=2pt]
\item downstream of separation the surface pressure no longer responds to the outer flow ---
it settles onto a nearly constant \emph{base pressure}, giving the flat plateau seen in
$\Cp(\theta)$ after the knee;
\item the wake is wide and low-pressure, so the rear of the cylinder no longer balances the
front: form drag appears, $\CD \approx 1.2$ for a laminar boundary layer.
\end{itemize}
Operationally, the separation angle is therefore read as \emph{the angle where the pressure
recovery ends and the plateau begins}, not as the point where $\Cp$ is lowest. In this
experiment the laminar separation point sits at roughly $\thsep = 80^\circ$--$90^\circ$ and
$\Cp^{\min} \approx -2.2$ rather than $-3$, because the boundary layer thickens and
displaces the effective body before it detaches.

\subsection{The trip wire, and why it is worth two extra runs}

A wire placed on the surface \emph{upstream of} the natural separation point forces the
boundary layer to transition to turbulence before it separates. A turbulent profile is fuller
(more momentum close to the wall), so it survives a much stronger adverse gradient: separation
moves aft to roughly $120^\circ$--$130^\circ$, the wake narrows, base pressure rises and
$\CD$ collapses from $\approx 1.2$ to $\approx 0.5$.

This is the same phenomenon as the drag crisis, reached at a far lower Reynolds number. The
whole point of running the same two speeds with and without the wire is that \CD{} is not a
function of $\ReD$ alone: at fixed $\ReD$ the boundary-layer \emph{state} changes \CD{} by
more than a factor of two. That contrast is the core result of the lab, and it is why the
second and third oral questions exist.

%=====================================================================
\section{The apparatus, and what each reading means}

\subsection{Tappings, tubes and reference heights}

Twenty-eight tappings at the angles tabulated in the handout (holes 1 and 7 at $\pm30^\circ$
mirrored about hole 2 at $0^\circ$) feed tubes 1--28. The Pitot-static tube is on tubes 35 and
36. Two consequences matter for the reduction:
\begin{itemize}[leftmargin=1.2em,itemsep=1pt,topsep=2pt]
\item hole 2 sits at the front stagnation point, so $\Cp(0^\circ)$ should come out at $+1$,
and its column height should equal the \emph{total} Pitot column;
\item holes 1 and 7 are mirror images, so an aligned cylinder must give
$\Cp(-30^\circ) = \Cp(+30^\circ)$. A mismatch is a rigid-body rotation of the whole ring and
is used to correct the datum (see Section~\ref{sec:align}).
\end{itemize}

\subsection{Which Pitot tube is static and which is total}

The static tube reads the smaller column (static pressure only); the total tube reads the
larger by exactly the dynamic head. Two independent checks identify them within the first
minute of the session:
\begin{enumerate}[leftmargin=1.4em,itemsep=1pt,topsep=2pt]
\item the tube whose column equals the hole-2 column is the \emph{total} tube;
\item $B = h_t - h_s > 0$ must hold, and $q = \rho_e g B$ must agree with
$\tfrac12 \rho U^2$ for the nominal tunnel speed ($\approx 146$\,Pa at $700$\,rpm,
$\approx 559$\,Pa at $1350$\,rpm).
\end{enumerate}
The code detects the assignment automatically and prints it; if it disagrees with what you
observe, set \texttt{cfg.pitotOrder} explicitly rather than silently accepting it.

\subsection{Alignment check}\label{sec:align}

With $g = \mathrm{d}\Cp/\mathrm{d}\theta$ evaluated at $+30^\circ$, a magnitude mismatch
$\Delta \Cp{}_{30} = \Cp(+30^\circ) - \Cp(-30^\circ)$ corresponds to a ring rotation
\begin{equation}
\Delta\theta_{\text{ring}} = -\frac{\Delta\Cp{}_{30}}{2g}.
\end{equation}
The factor 2 appears because both taps move in the same direction. Measured values of
$0.05^\circ$--$0.14^\circ$ are normal; a value above $1^\circ$ means the cylinder should be
rotated and the run repeated.

%=====================================================================
\section{Run plan}

\begin{table}[H]
\centering\small
\begin{tabularx}{\textwidth}{@{}llllX@{}}
\toprule
Run & Trip wire & Engine & Expected $U'$ & What it is for \\
\midrule
1 & none & 700 rpm & $\approx 15.5$\,m/s, $\ReD \approx 1.0\times10^{5}$ &
laminar separation, $\CD \approx 1.2$: the baseline \\
2 & none & 1350 rpm & $\approx 30.5$\,m/s, $\ReD \approx 2.1\times10^{5}$ &
same boundary-layer state, four times the dynamic head: cleaner data \\
3 & fitted & 700 rpm & $\approx 15.5$\,m/s &
tripped at the lower speed: separation far aft, low $\CD$ \\
4 & fitted & 1350 rpm & $\approx 30.5$\,m/s &
tripped at the higher speed: the pair used for the equivalent Reynolds number \\
5 & fitted & optional & --- &
only if the demonstrator offers a third speed; strengthens the $\CD$--$\ReD$ trend \\
\bottomrule
\end{tabularx}
\caption{Run plan. The two untripped runs and the two tripped runs form a
$2\times2$ design: two Reynolds numbers, two boundary-layer states.}
\end{table}

Between runs the tunnel is brought to the new speed and left to settle before any column is
read; the manometer needs tens of seconds and the ethanol menisci are slow. Read all 30
columns at one sitting per run, in one pass, so that a slow drift cannot masquerade as a
pressure distribution.

%=====================================================================
\section{Data reduction, in the order it has to be done}

\subsection{Step 1 --- \Cp{} for every tapping}

Equation~\eqref{eq:cp}. Nothing else can be computed before this, and the whole chain inherits
the reading error at this point.

\subsection{Step 2 --- uncorrected coefficients}

With the 27 angles from $0^\circ$ to $180^\circ$ (hole 1 at $-30^\circ$ is used only for the
symmetry check, not for the integral), the trapezoidal sums are
\begin{equation}
\CD' = \int_0^{\pi}\Cp\cos\theta\,\mathrm{d}\theta,\qquad
\CL' = \frac12\int_0^{2\pi}\Cp\sin\theta\,\mathrm{d}\theta .
\end{equation}
Only half the circumference is instrumented, so the missing half is supplied by symmetry
$\Cp(360^\circ-\theta) = \Cp(\theta)$; the antisymmetric part of the measured data is reported
separately as an asymmetry term ($\CL$ from the measured $-30^\circ$ point, essentially
$\Delta\Cp{}_{30}/2$). On a non-uniform grid the trapezoidal weights are not equal --- the
weights follow the actual angular spacing of each pair of taps.

\subsection{Step 3 --- blockage correction (Roshko, eqs 1--3)}

The cylinder blocks 22\,\% of the test-section breadth, so the flow around it is accelerated
relative to the nominal tunnel speed. With $r = d/b = 0.2217$ and $\CD'$ from step 2,
\begin{align}
\frac{V}{V'} &= 1 + \tfrac14 \CD'\,r + 0.82\,r^2, \label{eq:vv}\\
\CD{}_{\text{eq2}} &= \left(1 - \tfrac12 \CD'\, r - 2.5\,r^2\right)\CD', \label{eq:cd2}\\
\Cp &= 1 + \left(\frac{V'}{V}\right)^{2}\left(\Cp' - 1\right). \label{eq:cp3}
\end{align}
The order is forced: \eqref{eq:vv} needs $\CD'$, and \eqref{eq:cp3} needs $V/V'$. The corrected
velocity is $V = V'\,(V/V')$, and the corrected Reynolds number follows from it.

\subsection{Step 4 --- corrected drag, two ways}

\begin{itemize}[leftmargin=1.2em,itemsep=1pt,topsep=2pt]
\item integrate the corrected distribution from \eqref{eq:cp3}: $\CD{}_{\text{int}}$;
\item apply the direct correction \eqref{eq:cd2} to $\CD'$: $\CD{}_{\text{eq2}}$.
\end{itemize}
These are not the same number and should not be forced to agree: \eqref{eq:cd2} is an
empirical fit to the whole flow, while $\CD{}_{\text{int}}$ only knows the 28 sampled pressures.
The gap is typically 7--9\,\% and is itself a result worth commenting on.

\subsection{Step 5 --- equivalent Reynolds number for the tripped runs}

$\CD$ alone does not fix $\ReD$, because the $\CD(\ReD)$ curve is double-valued around the
critical point. Map the tripped $\CD$ onto the \emph{rising} branch of the reference curve
(Roshko Fig.~2 / Anderson Fig.~3.44) and read off $\ReD_{\text{eq}}$. Expect
$\ReD_{\text{eq}} \sim 10^6$ from the tripped runs, one to two decades above the actual
tunnel $\ReD$ --- which is exactly the point of tripping: it buys a high-\ReD{} boundary layer
at low-\ReD{} cost.

\subsection{Step 6 --- separation angle}

From the \Cp{} distribution: locate the suction peak, then walk aft until the recovery slope
$\mathrm{d}\Cp/\mathrm{d}\theta$ stays small over a 12$^\circ$ window. The criterion, the
thresholds and their justification are in the companion document; the operational definition
is ``the angle where the plateau begins'', and for a smeared knee the answer carries
$\pm 20^\circ$ rather than $\pm 2.5^\circ$.

%=====================================================================
\section{Uncertainty plan}

The full derivation is in the companion notes; what has to be decided \emph{before} the lab is
which error dominates and therefore where to spend effort.

\begin{table}[H]
\centering\small
\begin{tabularx}{\textwidth}{@{}lXl@{}}
\toprule
Quantity & How it enters & Expected $1\sigma$ \\
\midrule
One column height $h_i$ & type A from 3 repeats $+\ \pm0.5$\,mm scale term; $\Cp$ error $=u(h)/B$ &
0.027 (low speed), 0.007 (high) \\
Tap angle $\theta_i$ & machining tolerance $0.5^\circ/\sqrt3$, per tap & $0.29^\circ$ \\
Ring rotation & from the $\pm30^\circ$ pair, coherent over all taps & $0.05$--$0.14^\circ$ \\
Air temperature & type B only ($0.5$\,K half-width), enters $\rho$ and $\mu$ & $0.3$\,K \\
Barometric pressure & type B only, enters $\rho$ & 50\,Pa \\
Separation angle & slope noise divided by knee curvature, floored at half the tap spacing & $2.5^\circ$ or $\gg$ \\
\bottomrule
\end{tabularx}
\caption{Where the uncertainty comes from. The low-speed runs are four times noisier in \Cp{}
than the high-speed ones, purely because $B$ is four times smaller.}
\end{table}

Two practical consequences to act on at the bench:
\begin{enumerate}[leftmargin=1.4em,itemsep=1pt,topsep=2pt]
\item \textbf{read each tapping at least three times}; with one reading the type-A term
disappears and $u(\CD')$ roughly triples;
\item \textbf{prefer the high-speed runs} when a single number has to be quoted, and say
so explicitly --- $u(\CD')$ is 1.0\,\% against 4.0\,\% at the low speed.
\end{enumerate}

Report every $u$ as $1\sigma$, and state $k=2$ where a 95\,\% figure is quoted. Do not mix the
two conventions on the same graph.

%=====================================================================
\section{Checks to run at the bench}

\begin{table}[H]
\centering\small
\begin{tabularx}{\textwidth}{@{}lX@{}}
\toprule
Check & Expected \\
\midrule
$\Cp(0^\circ)$ from hole 2 & $+1.00 \pm 0.05$, and equal to the total Pitot column \\
$\Cp(90^\circ)$ untripped & $-2.0$ to $-2.3$ (ideal $-3$) \\
$\Cp(90^\circ)$ tripped & $-1.4$ to $-1.6$ \\
$\Cp(\pm30^\circ)$ mismatch & below $0.01$ in $\Cp$, i.e.\ $\Delta\theta_{\text{ring}} < 0.2^\circ$ \\
Base pressure $180^\circ$ & $\Cp \approx -0.3$ untripped, $\Cp \approx -0.1$ to $-0.2$ tripped \\
$\CD'$ untripped / tripped & $\approx 1.1$ / $\approx 0.5$ \\
Separation, untripped / tripped & $85^\circ$--$95^\circ$ / $120^\circ$--$135^\circ$ \\
\bottomrule
\end{tabularx}
\caption{If any of these fails by more than the tolerance, suspect a tube leak, a bubble or a
mistaken Pitot assignment before suspecting the theory.}
\end{table}

%=====================================================================
\section{Oral assessment: the eleven questions, with answers}\label{sec:oral}

The handout lists the questions below. They are reproduced with the answers I intend to give,
so that the numbers and the phrasing are settled beforehand.

\begin{enumerate}[leftmargin=1.6em,itemsep=4pt,topsep=2pt]

\item \textbf{How do we vary Reynolds number, and why this way? Which speed is more reliable?}\\
$\ReD = \rho U d/\mu$ and $d$, $\rho$, $\mu$ are fixed by hardware and weather, so the only
lever is $U$, i.e.\ engine rpm. The high speed is more reliable: $q \propto U^2$, so the
Pitot column doubles while the reading error stays 0.5\,mm, and the relative \Cp{} error
falls from 1.9\,\% to 0.5\,\%.

\item \textbf{Where is the best place to put the wire?}\\
On the surface, upstream of the natural separation point --- about $60^\circ$--$75^\circ$
from the front stagnation point. Too far forward and transition happens too early, giving a
over-thick turbulent layer and an unrepresentative \Cp{}; too far aft (past separation) and
it does nothing, because the layer has already detached. Diameter of order 1\,mm, spanning
the full width at mid-span.

\item \textbf{Why does separation differ with and without the wire at the same \ReD?}\\
Because separation is set by the boundary-layer \emph{state}, not by $\ReD$ alone. A laminar
profile cannot sustain much adverse gradient and detaches near $85^\circ$; a turbulent profile
carries more momentum near the wall and survives to $130^\circ$. Same $\ReD$, different
flow: this is the clearest demonstration that $\CD$ is not a single-valued function of $\ReD$.

\item \textbf{What effect does this have on drag?}\\
Delaying separation raises the base pressure (the wake is narrower and less suctorial), so
form drag collapses: $\CD$ from $\approx 1.2$ to $\approx 0.5$. Skin friction rises slightly
with turbulence, but the form-drag drop dominates by an order of magnitude.

\item \textbf{What is the effective Reynolds number for the tripped case?}\\
Take the measured tripped $\CD$ (about 0.48--0.50) and follow the rising branch of the
$\CD(\ReD)$ curve until it reaches that value: $\ReD_{\text{eq}} \approx 1$--$2\times10^{6}$,
one to two decades above the tunnel $\ReD \approx 2\times10^{5}$.

\item \textbf{Two examples where tripping a boundary layer is advantageous.}\\
A golf ball: dimples force early transition, delay separation and roughly double the carry.
A turbine or compressor blade at low Reynolds number: a trip strip on the suction side
prevents laminar separation, which would otherwise stall the blade --- the same mechanism
used on low-\ReD{} wind-tunnel models so that they reproduce full-scale behaviour.

\item \textbf{What does blockage do to the pressure measurements and to drag?}\\
It accelerates the flow around the model relative to the nominal tunnel speed and constrains
the wake. Measured suction is too strong and measured $\CD$ too high: the sampled
$\Cp^{\min} \approx -2.27$ is corrected to $\approx -1.70$, and the direct correction reduces
$\CD'$ by 5--10\,\%. The correction factor here is $V/V' = 1.102$ at the low speed.

\item \textbf{Other sources of error besides blockage?}\\
Column reading (parallax, meniscus) and scale resolution; tube leaks, trapped bubbles, clogged
taps; thermal drift over the session; unsteady vortex shedding, which makes a single reading
a sample of a fluctuating field rather than a mean; three-dimensional end effects and wall
boundary layers; ring misalignment; surface roughness and any wire actually changing the
geometry; and the finite number of taps, which under-resolves the suction peak.

\item \textbf{Which of the two Pitot readings is static and which is total?}\\
The smaller column is static; the larger is total. Confirmed at the bench by matching the
total tube to hole 2 (stagnation), and by checking $q = \rho_e g B$ against
$\tfrac12\rho U^2$ for the nominal rpm.

\item \textbf{What speed would give a turbulent boundary layer without a trip wire?}\\
Natural transition on a smooth cylinder begins around $\ReD \approx 2\times10^{5}$ and is
complete by $\approx 3.5\times10^{5}$, i.e.\ $U = \ReD\mu/(\rho d) \approx 30$--$50$\,m/s.
The upper end is beyond this tunnel's practical range (and the model would vibrate), which
is precisely why the wire is needed.

\item \textbf{What other surface features could force transition?}\\
Distributed roughness (grit or sandpaper); dimples; spanwise trip dots or a strip of
tape/thickness change; a seam or step; or a small spanwise protrusion. All work by
introducing a local disturbance whose length scale is comparable to the boundary-layer
thickness.

\end{enumerate}

%=====================================================================
\section{Logistics and what to write in the logbook}

\begin{itemize}[leftmargin=1.2em,itemsep=2pt,topsep=2pt]
\item Arrive before the stated start; late entry is refused after two minutes.
\item Note the tunnel air temperature (three readings) and barometric pressure (three readings)
\emph{before} the first run, and the manometer scale increment (mm per division).
\item Fill one row per tapping per run in \texttt{readings.csv} (wide layout: three repeat
columns), and record which Pitot tube was static.
\item After each run, note the engine rpm, whether the wire was fitted and where, and any
anomaly (bubble, leaking tube, drifting zero).
\item In the logbook, record the seven numbers per run that matter: $\CD'$,
$\CD{}_{\text{eq2}}$, $\CD{}_{\text{int}}$, $V/V'$, $\thsep$, $\ReD'$, plus
$u(\CD')$, $u(\thsep)$ and $\Delta\theta_{\text{ring}}$.
\item Keep the raw column heights. Everything else can be recomputed from them; they cannot.
\end{itemize}

\end{document}
